% 原书分章 PDF 第 11—14 页：第 22.3.3 节与第 22.4 节。
\subsection{真实最优规则\texorpdfstring{$\psi_0$}{ψ₀}的统计推断}

现在假定我们关心的是最优规则 $d_0$ 下平均结局的估计，而非数据自适应参数 $\tilde\psi_{0n}$。注意，
\[
\begin{aligned}
 \sqrt n(\tilde\psi_n^*-\psi_0)
 &=\sqrt n(\tilde\psi_n^*-\tilde\psi_{0n})
 +\sqrt n(\tilde\psi_{0n}-\psi_0)\\
 &=\sqrt n(\tilde\psi_n^*-\tilde\psi_{0n})
 +\frac{\sqrt n}{J}\sum_{j=1}^J[\Psi_{d_{nj}}(P_0)-\psi_0].
\end{aligned}
\]
如果 $J^{-1}\sum_{j=1}^J[\Psi_{d_{nj}}(P_0)-\psi_0]$，那么由 Slutsky 定理，只要定理 22.4 的条件成立，左侧就具有与 $\sqrt n(\tilde\psi_n^*-\tilde\psi_{0n})$ 相同的正态极限。此外，由于 $n\to\infty$ 时 $J$ 固定，若
\begin{equation}
 \Psi_{d_{nj}}(P_0)-\psi_0=o_P(n^{-1/2})\qquad\text{对每个 }j\text{ 都成立},
\tag{22.8}
\end{equation}
则 $J^{-1}\sum_{j=1}^J[\Psi_{d_{nj}}(P_0)-\psi_0]=o_P(n^{-1/2})$。

为分析 $\Psi_{d_{nj}}(P_0)-\psi_0$，我们将假定用户用 $\bar Q_{b,1nj}$ 和 $\bar Q_{b,2nj}$ 估计 $\bar Q_{b,10}$ 和 $\bar Q_{b,20}$，然后使用定理 22.1 所述形式的代入估计量（plug-in estimator）。Luedtke 和 van der Laan（2016b）此前已介绍 $\bar Q_{b,10}$ 与 $\bar Q_{b,20}$ 的数据自适应估计量。虽然我们并不要求 $d_{nj}$ 必须来自代入估计量，但这里将集中分析这种估计方案。对于代入估计量，证明式 (22.8) 所需的主要结果本质上是一个解析结果，因此我们先考虑一般的 $Q$，以及相应的处理增益函数 $\bar Q_{b,1},\bar Q_{b,2}$ 和最优规则的代入估计 $d_{Q,A(0)},d_{Q,A(1)}$。然后就可以把这一结果直接应用于我们的各折特定估计量。

下列结果的证明见第 22.5 节。

\begin{bookresult}{引理}{22.1}
回顾定理 22.1 中 $\bar Q_{b,20}$ 和 $\bar Q_{b,10}$ 的定义。$\Psi(P_0)=\E_0Y_{d_0}$ 可以表示为
\[
\begin{aligned}
 \Psi(P_0)
 &=\E_0Y_{(0,1),(0,1)}\\
 &\quad+\E_0\bigl[d_{0,A(1)}((0,1),V_{(0,1)}(1))
 \bar Q_{b,20}((0,1),V_{(0,1)}(1))\bigr]\\
 &\quad+\E_0d_{0,A(0)}(V(0))\bar Q_{b,10}(V(0)).
\end{aligned}
\]
其中，$V_{(0,1)}(1)$ 从第一时间点给予处理 $(0,1)$ 的 $g$ 计算分布中抽取。
\end{bookresult}

由此可得
\[
\begin{aligned}
 R_2(Q,Q_0)
 &=\E_0(d_{Q,A(0)}-d_{0,A(0)})(V(0))\bar Q_{b,10}(V(0))\\
 &\quad+\E_0(d_{Q,A(1)}-d_{0,A(1)})((0,1),V_{(0,1)}(1))
 \bar Q_{b,20}((0,1),V_{(0,1)}(1))\\
 &\equiv R_{2,A(0)}(Q,Q_0)+R_{2,A(1)}(Q,Q_0).
\end{aligned}
\]
在间隔假设（margin assumption）下，我们能够使 $R_2(Q,Q_0)$ 达到较快的收敛率。先给出用于界定 $R_{2,A(0)}(Q,Q_0)$ 的假设。假定存在正常数 $C_1,\beta_1$，使得对所有 $t>0$，
\begin{equation}
 P_0\bigl\{0<|\bar Q_{b,10}(V(0))|\leq t\bigr\}\leq C_1t^{\beta_1}.
\tag{MA1}
\end{equation}
上述假设要求，第一时间点的处理增益函数不要在决策边界（零）附近（但不包括边界本身）聚集过多概率质量。这个假设不同于例外分布条件，因为后者要求该处理增益函数恰好在决策边界上不放置任何概率质量。$\beta_1$ 和 $\beta_2$ 较小时，这是一个较弱的假设，但此时可能达不到满足式 (22.8) 所需的收敛率。如果把处理增益函数应用于数据所得的随机变量，在零附近具有有界的 Lebesgue 密度，那么 $\beta_1=1$ 时这些假设成立。

我们对 $\bar Q_{b,20}$ 作类似假设。具体而言，假定存在某些 $C_2,\beta_2$，使得对所有 $t>0$，
\begin{equation}
 P_0\bigl\{0<|\bar Q_{b,20}((0,1),V_{(0,1)}(1))|\leq t\bigr\}
 \leq C_2t^{\beta_2}.
\tag{MA2}
\end{equation}
现在证明，(MA1) 与 (MA2) 可以利用 $\bar Q_{b,1}$、$\bar Q_{b,2}$ 到 $\bar Q_{b,10}$、$\bar Q_{b,20}$ 的距离，为 $R_2(Q,Q_0)$ 给出依赖 $\beta_1,\beta_2$ 的上界。

\begin{bookresult}{定理}{22.5}
如果对某些 $C_1,\beta_1>0$，(MA1) 成立，那么存在某个常数 $C>0$，使得
\begin{equation}
\begin{aligned}
 |R_{2,A(0)}(Q,Q_0)|\leq C\min\Bigl\{
 &\|\bar Q_{b,1}-\bar Q_{b,10}\|_{2,P_0}^{2(1+\beta_1)/(2+\beta_1)},\\
 &\|\bar Q_{b,1}-\bar Q_{b,10}\|_{\infty,P_0}^{1+\beta_1}\Bigr\}.
\end{aligned}
\tag{22.9}
\end{equation}
如果对某些 $C_2,\beta_2>0$，(MA2) 成立，那么存在某个常数 $C>0$，使得
\begin{equation}
\begin{aligned}
 |R_{2,A(1)}(Q,Q_0)|\leq C\min\Bigl\{
 &\|\bar Q_{b,2}-\bar Q_{b,20}\|_{2,P_{0,(0,1)}}^{2(1+\beta_1)/(2+\beta_1)},\\
 &\|\bar Q_{b,1}-\bar Q_{b,10}\|_{P_{0,(0,1)}}^{1+\beta_1}\Bigr\},
\end{aligned}
\tag{22.10}
\end{equation}
其中，$P_{0,(0,1)}$ 表示在第一时间点给予处理 $(0,1)$ 这一静态干预特定的 $g$ 计算分布。
\end{bookresult}

当 $d_{nj}$ 用代入估计量估计时，上述解析结果有助于判断式 (22.8) 是否合理。在参数模型中，通常可以在上述 $L_2(P_0)$ 范数和上确界范数下，以 $n^{-1/2}$ 的速率估计 $\bar Q_{b,10}$ 和 $\bar Q_{b,20}$。因此，如果间隔条件（margin condition）在 $\beta_1=\beta_2=1$ 时成立，上确界范数的结果就使每个 $\Psi_{d_{nj}}(P_0)-\psi_0$ 达到 $n^{-1}$ 的速率。实践中，我们当然不指望能够正确指定一个参数模型；相反，会使用超级学习等数据自适应估计量来估计处理增益函数，以提高估计量被正确设定的可能性（Luedtke 和 van der Laan，2016b）。在对处理增益函数作出光滑性假设后，可以利用平滑方法，保证在 $L_2(P_0)$ 范数下达到接近参数模型的速率。如果足够光滑，这些速率能够达到上述定理在 $\beta_1=\beta_2=1$ 时要求的 $o_P(n^{-3/8}$ 速率，从而说明 $\Psi_{d_{nj}}(P_0)-\psi_0=o_P(n^{-1/2})$。然而，一般而言，我们未必能期待如此快的速率。如果这一快速率不成立，仍然可以对数据自适应参数 $\tilde\psi_{0n}$ 进行推断。如果快速率成立（当 $V$ 维数较低时可能如此），那么在定理 22.4 的条件及 $\sigma_n^2$ 的一致性下，式 (22.7) 给出的置信区间对 $\tilde\psi_{0n}$ 和 $\psi_0$ 都渐近有效。

\section{讨论}

本章研究了 $V$ 最优规则下平均结局的半参数统计推断，以及一个数据自适应目标参数的统计推断。后一个参数定义为数据自适应确定的 $V$ 最优规则下的平均结局，并将这一规则视为给定。我们证明了一个令人意外且有用的结果：当数据由非例外分布生成时，$V$ 最优规则下的平均结局可表示为一个统计参数，其路径导数与把未知规则视为已知时完全相同。因此，可以立即得到有效影响曲线；针对给定规则下平均结局的任何有效估计量，都可以应用于估计得到的规则。具体而言，我们给出了一个 CV-TMLE，并给出渐近线性结果。不过，统计目标参数对未知规则的依赖会影响 CV-TMLE 的二阶项。因此，CV-TMLE 的渐近线性要求估计规则与 $V$ 最优规则之间的某个二阶差异，以快于 $1/\sqrt n$ 的速率收敛于零。对于仅依赖一个连续评分（如某个生物标志物）的规则，可以期待这一要求成立；但当 $V$ 为中高维时，只有很强的光滑性假设才能保证这一点。因此，即使在随机对照试验中，也不能期待对这样的 $V$ 最优规则作出有效统计推断。

为应对这一挑战，我们还介绍了各样本切分特定的数据自适应目标参数之平均的估计，其一般思路由 Hubbard 等（2016）提出。具体而言，对 $J$ 次训练样本与验证样本的切分，分别考虑在训练样本上拟合的动态处理下的平均结局，再取这 $J$ 个量的平均，以此定义数据自适应目标参数。我们给出了这一数据自适应目标参数的 CV-TMLE，并建立渐近线性定理；该定理不要求估计规则一致收敛于最优规则，更不要求以某个特定速率收敛。我们说明，对这个数据自适应目标参数的统计推断不依赖估计规则向最优规则的收敛率；事实上，只要求数据自适应拟合的规则收敛于某个固定的、可能次优的规则。因此，在序贯随机对照试验中，这个方法在很温和的条件下就能给出有效的渐近统计推断；其中最主要的条件，是估计规则收敛于某个固定的、可能次优的规则。

对最优处理策略作出推断是一个重要主题，有望帮助指导未来的卫生政策决策。我们认为，使用一个较大的半参数模型是可取的，因为这有助于保证：预期实施某个估计处理策略所带来的健康收益，并非源于模型错误设定造成的偏差。本章提出的 CV-TMLE 具有许多理想的统计性质，使我们能够在这个大模型中进行估计与推断。
